\chapter{Methods}
\label{ch-2}

This thesis contains ten chapters. The first two chapters introduce the research topic and outline the methods used. The middle seven chapters form the substantive part of the thesis and generate new knowledge. The final chapter discusses the findings of the research and concludes the thesis.

Chapter \ref{ch-3} aimed to describe and demonstrate an approach to include environmental impacts in a cost analysis. It used standard economic methods to understand the feasibility of doing so applied to the pilot MEL-SELF RCT. 

Chapter \ref{ch-4} aimed to review methods to include environmental impacts in health economic evaluations and health technology assessments. It used methods to conduct a scoping review recommended by JBI (formerly an acronym for the Joanna Briggs Institute).

Chapter \ref{ch-5} aimed to estimate the health system costs and associated carbon emissions of diagnostic imaging tests used by patients undergoing different surveillance schedules for Stage III melanoma follow up. It used epidemiological and economic methods to estimate the costs and carbon emissions of a defined patient cohort. The results of this chapter were used in the cost-effectiveness analysis reported in Chapter \ref{ch-7}.

Chapter \ref{ch-6} aimed to estimate the carbon emissions of appointments at a skin cancer clinic related to the identification and excision of lesions suspicious for melanoma and other skin cancer. It used process-based life cycle assessment methods to estimate carbon emissions of each type of appointment. The results of this chapter were used in the cost-effectiveness analysis reported in Chapter \ref{ch-7}.

Chapter \ref{ch-7} aimed to demonstrate methods to include carbon emissions in a cost-effectiveness analysis. It used standard economic methods for the base case analysis. Methods identified in Chapter \ref{ch-4} were used to include environmental impacts, and these were evaluated. 

Chapter \ref{ch-8} aimed to describe items to include in a health economic analysis plan (HEAP) to prespecify methods to include environmental impacts. The HEAP for the cost consequence analysis of the MEL-SELF trial was used as a case study. 

Chapter \ref{ch-9} aimed to outline the current Australian policy context for including environmental impacts in HTA, evaluate potential methods to include environmental impacts for public reimbursement, and detail a potential method for the cost-minimisation pathway. 

Chapter \ref{ch-10} summarises key findings, explores their implications for researchers and policy makers, and concludes the thesis.